\section{有限要素法：場を節点値で近似する}
\label{theory:fem:background}
連続体の変位を有限個の節点値で表し、材料則とつり合いを離散方程式へ移す。変位を$u_h=\sum_iN_i u_i$、微分演算を$L$とすると、ひずみは$\varepsilon=Lu_h=Bq$、線形弾性の応力は$\sigma=D\varepsilon$である。形状関数の微分を含む$B$を用い、要素剛性は$K_e=\int_{V_e}B^TDB\,dV$となる。この補間から剛性行列への関係はMIT講義でも整理される\parencite{mit_roy_fea1999}。以下の具体的な形状関数と幾何計算は独自に導出する。

\subsection{一次三角形：一定ひずみの理由}
\label{theory:fem:triangle}
非退化の三頂点を持つ一次三角形では、各変位成分を$u=a+bx+cy$と表す。微分は$u_x=b,u_y=c$であるため、微小ひずみは要素内で一定となる。均質な線形弾性なら応力も一定である。6節点の二次三角形には$x^2,xy,y^2$の項が加わり、微分は一般に一次関数となる。

これは任意の節点変位に対する要素の性質である。二次三角形でも二次項がゼロとなる特別な場ならひずみは一定になり得る。「一定ひずみ要素」を選ぶ設問で、その特殊例だけを理由に高次要素まで選ばない。退化三角形や不均質な材料では、一定ひずみから一定応力への移行にも条件が要る。

\subsection{双一次四角形：交差項が残る}
\label{theory:fem:q4}
自然座標の四隅を$(-1,-1),(1,-1),(1,1),(-1,1)$とすると
\[
 (N_1,N_2,N_3,N_4)=\frac14\bigl((1-\xi)(1-\eta),(1+\xi)(1-\eta),
 (1+\xi)(1+\eta),(1-\xi)(1+\eta)\bigr).
\]
各節点で対応する関数だけ1となり、和は1である。展開した$\xi,\eta,\xi\eta$の係数ベクトルは順に
\[
 \tfrac14(-1,1,1,-1),\quad\tfrac14(-1,-1,1,1),\quad
 \tfrac14(1,-1,1,-1).
\]
非定数の各係数の和が0という検算もできる。

長方形の物理座標なら$u=a+bx+cy+dxy$となり、$u_x=b+dy,u_y=c+dx$には位置依存が残る。一般四角形では座標写像のJacobianを介して微分し、ひずみはさらに形状に依存する。四隅4節点というだけで要素内一定応力と判断しない。

\subsection{変位の連続と応力の跳び}
\label{theory:fem:continuity}
適合する変位法では隣接要素が同じ辺の節点変位を共有する。一次の辺補間が同じなので辺上の変位も一致する。しかし法線方向の微分までは一般に一致せず、微分から作るひずみや応力に跳びが生じ得る。一様ひずみの場では隣接応力が一致する例もあるため、「常に不連続」という全称は強すぎる。

表示用に節点応力を平均すると滑らかな色分布になることがあるが、平均処理が近似精度の保証を与えるわけではない。変位、積分点応力、平均応力のどれを比較したかを明示する。

\subsection{メッシュ：局所応力に上下保証はない}
\label{theory:fem:mesh}
細分化するh改良は空間解像度を増し、高次化するp改良は近似関数の表現力を増す。滑らかな解、健全な形状、適切な積分や境界条件の下で有効となる。高勾配域へ要素を配置する考え方はMITの例とも対応する\parencite{mit_roy_fea1999}。

線形弾性の適合変位法における全体のエネルギーの傾向から、任意の点の応力が上界であるとはいえない。粗いメッシュでは応力集中を過小評価することもある。「変形が小さいので応力は常に高く安全側」という結論は成立しない。異なる荷重では重要な勾配の位置も変わり、一つの荷重で妥当だった分割を無検証で使い回さない。

複数のメッシュで同じ変位、反力、位置の応力を比較する。尖った角や点荷重に理想化した特異点ではピーク応力が細分ごとに増す場合があり、有限値への収束を当然視しない。要素の歪みや退化、丸め誤差、材料モデルの誤差は別に確認する。

\subsection{面積座標：重心・内心・外心を分ける}
\label{theory:fem:area}
三角形ABC内の点Pに対し、$S_A=S_{PBC},S_B=S_{PCA},S_C=S_{PAB}$とおき、$\lambda_i=S_i/S$と定める。非負で和は1となる。頂点Aでは$(1,0,0)$、重心では$(1/3,1/3,1/3)$である。面積座標は一次三角形の形状関数にもなる。

内心は三辺からの距離$r$が等しく、$S_A=r|BC|/2$などとなる。従って座標は向かい辺の長さの比である。$AB:BC:CA=3:4:5$なら$(\lambda_A,\lambda_B,\lambda_C)=(4,5,3)/12$となる。Bが直角で、外心は斜辺ACの中点なので$(1/2,0,1/2)$である。重心を内心や外心と混同せず、図の見た目より辺長比を優先する。三角形外の点に拡張する場合には符号付き面積が必要となる。

\subsection{二層の熱伝導：並列と直列の加算量}
\label{theory:fem:thermal}
定常伝導の平板では$Q=kA\Delta T/d$であり、熱流は温度勾配に比例する\parencite{openstax_thermal_conduction}。接触熱抵抗がなく、厚さ$l_1,l_2$の層に平行な流れでは温度勾配が共通であり、各断面の熱流を足す。層を横切る流れでは熱流が共通であり、各層の温度差を足す。$L=l_1+l_2$として
\[
 k_\parallel=\frac{k_1l_1+k_2l_2}{L},\qquad
 k_\perp=\frac{L}{l_1/k_1+l_2/k_2}.
\]
前者は厚さで重み付けした算術平均、後者は調和平均である。正の熱伝導率では両者とも各材料値の間に入り、$k_\parallel\ge k_\perp$となる。界面熱抵抗や過渡現象があれば、この簡単な二式のモデルをそのまま適用しない。

厚さを材料に逆に対応させた誤答式は等厚の特殊値で正しい式と一致する。試験では一般式を導く目的と読みつつ、その一致と時間条件の省略を解説の疑義へ残す。
